\section*{Figure captions}
% generated by make_submission.py; do not edit by hand
\noindent\textbf{Figure 1.} How the dataset was built. $n$ counts scans, or corrections, at each step.

\noindent\textbf{Figure 2.} Four phenotypes at the lumbar borders, aligned on the sacrum. Red, lowest rib-bearing vertebra and its rib; blue, sacrum; yellow, the lumbar bodies. Cases 0704, 0428, 0094, 0005, each the most ordinary member of its phenotype by rule. (b) and (c) present the same upper border, a short rib pair under a full pair, and were named differently.

\noindent\textbf{Figure 3.} Morphometry by level, T11--L5: median, interquartile range, 5th--95th percentile, $n$ at right. (a) End-plate width EPWu. (b) Canal width SCW and depth SCD. (c) Pedicle width PDW, both sides averaged. L5 is withheld in (a) and (c); see text. Reference is Panjabi's cadaveric series\cite{panjabi1991,panjabi1992}, keyed in (a), its SD recovered from the reported SEM; T11--T12 EPWu is digitised from his Figure 4A.

